% !TEX root = ../main.tex
\section{EndorseRank and AWP Side by Side}
\label{ch:appendix-er-awp}

Chapter~\ref{ch:results} reports every score of the thesis together and reads each PageRank against the raw degree it smooths. This appendix collects the direct comparisons between the two single-layer scores, EndorseRank and AWP: their family means and paired differences in the same window, their per-proxy coefficients, and their paired differences out of window. Apart from the trader grid at its end, none of these comparisons was fixed in advance, and none carries a claim of the thesis. The two scores walk on different edges and are built for different constructs, so on each check the score ahead is, almost every time, the one whose edges the check counts. Outcomes built from neither edge type are the exception, since neither score was built for them, and are reported in Section~\ref{sub:neutral-label-results}. Their coverage differs as well. No allowance edge touches $523$ of the matched contracts and no transfer edge touches $2{,}670$, and $2{,}891$ of the $17{,}050$ spenders of W0 have no transfer edge before the freeze. Out of window AWP ranks later new senders above the transfer in-degree in both windows, while EndorseRank ranks later new approval pairs below the in-approve degree in both (Section~\ref{sec:holdout-results}), so a lead of EndorseRank over AWP does not show that propagation over approvals adds anything to the count of approving owners.

\dissertationSubheading[sec:er-awp-same-window]{Same window}

Table~\ref{tab:alignment-hybrid} gives the family means of the two scores with their intervals, and the last row of Table~\ref{tab:tie-sensitivity} their agreement. The two patterns cross. AWP tracks the transfer family most closely ($0.572$, interval $[0.564,0.580]$) and the allowance family least ($0.193$, $[0.182,0.203]$), while EndorseRank tracks the allowance family most closely ($0.585$, $[0.578,0.592]$) and reaches $0.158$ ($[0.147,0.169]$) on the transfer family. Table~\ref{tab:tau-diff} tests five paired differences on the same contract resamples (Section~\ref{sec:statistical-tests}), numbered i to v in the order of its rows.

% Generated by experiments/scripts/export_latex.py -- do not edit by hand.
% Source: experiments/data/2-processed-tables-and-evaluations/revised-usd-12m/same-window/eval_summary.json
\begin{table}[htbp]
\centering
\caption[Same-window paired contrasts between EndorseRank and AWP.]{Paired bootstrap tests of differences between Kendall $\tau$ coefficients of EndorseRank and AWP in the same window (matched cohort, $n=6{,}483$ contracts). Family entries use the family-mean $\tau$; the same contract resamples are used for $a$ and $b$, so the interval is on the paired difference. An interval excluding 0 indicates a difference not attributable to sampling variation at the 5\% level. CI = 95\% percentile bootstrap over contracts (400 paired resamples, seed 42).}
\label{tab:tau-diff}
\fitwidth{%
\begin{tabular}{lccccc}
\toprule
Contrast ($a$ minus $b$) & $\tau(a)$ & $\tau(b)$ & $\Delta\tau$ & 95\% CI of $\Delta\tau$ & $P(\Delta\tau>0)$ \\
\midrule
EndorseRank allowance minus EndorseRank transfer & 0.560 & 0.196 & +0.363 & [0.341, 0.382] & 1.000 \\
AWP transfer minus EndorseRank transfer & 0.736 & 0.196 & +0.539 & [0.518, 0.560] & 1.000 \\
EndorseRank allowance minus AWP allowance & 0.560 & 0.226 & +0.333 & [0.311, 0.354] & 1.000 \\
AWP sybil-stability minus EndorseRank sybil-stability & 0.419 & 0.026 & +0.393 & [0.373, 0.412] & 1.000 \\
EndorseRank in-approve degree minus EndorseRank in-degree & 0.650 & 0.218 & +0.432 & [0.410, 0.452] & 1.000 \\
\bottomrule
\end{tabular}
}
\end{table}


All five intervals lie above zero. AWP's transfer alignment is higher than EndorseRank's (contrast~ii, $+0.414$, $[0.401,0.427]$), EndorseRank's allowance alignment is higher than AWP's (contrast~iii, $+0.392$, $[0.380,0.405]$), and AWP leads EndorseRank on the Sybil-stability family (contrast~iv, $+0.388$, $[0.376,0.401]$). The remaining two compare EndorseRank with itself: its allowance alignment is higher than its transfer alignment (contrast~i, $+0.427$, $[0.416,0.440]$), and its in-approve-degree alignment is higher than its in-degree alignment (contrast~v, $+0.473$, $[0.461,0.486]$). EndorseRank ties only $0.17\%$ of the contract pairs, so its own ties cap none of these coefficients, and the ties of the proxies do not explain gaps of this size: on the in-degree AWP reaches $0.578$ ($[0.569,0.587]$), where EndorseRank reaches $0.182$ ($[0.169,0.194]$). With the contracts outside the allowance graph kept as isolated nodes instead of being scored zero, none of the family means changes (Table~\ref{tab:tie-sensitivity}).

The opposite directions of contrasts~ii and~iii follow from the constructs. Each score is ahead on the family built from its own edges, and the main proxies of each family include the in-degree that the corresponding PageRank smooths, so the agreement is partly built in and checks the intended construct of each score. Neither difference says which score is the better measure of reputation.

On the scaling stages of Table~\ref{tab:robustness-sample-size}, whose values are not comparable with those of the full cohort, the order of the two scores within each family is the same at every stage: AWP is ahead on the transfer and Sybil-stability families and EndorseRank on the allowance family.

\dissertationSubheading[sec:proxy-family-detail]{Proxy detail}

Each family mean condenses several per-proxy coefficients. Tables~\ref{tab:alignment-transfer} to~\ref{tab:alignment-sybil-stability} report them one proxy at a time for EndorseRank and AWP, as Spearman $\rho$, Kendall $\tau$ and the 95\% interval of $\tau$. Before correlation, every proxy is oriented so that a higher value means higher expected reputation, and the transfer proxies are the ones used in the evaluation of AWP by its authors \parencite{do2023awp}.

\dissertationSubsubheading[sub:transfer-family]{Transfer family}

Table~\ref{tab:alignment-transfer} gives the coefficients of the two scores with the in-degree and the in-value.

% Generated by experiments/scripts/export_latex.py -- do not edit by hand.
% Source: experiments/data/2-processed-tables-and-evaluations/revised-usd-12m/same-window/eval_summary.json
\begin{table}[htbp]
\centering
\caption[Domain alignment of EndorseRank and AWP with the transfer proxies.]{Domain alignment of EndorseRank and AWP with the transfer proxies (in-degree and in-value of each contract), same window (matched cohort, $n=6{,}483$ contracts). Kendall $\tau$ and Spearman $\rho$ between scores and proxy rankings. CI = 95\% percentile bootstrap over contracts (400 paired resamples, seed 42).}
\label{tab:alignment-transfer}
\fitwidth{%
\begin{tabular}{lccc}
\toprule
Proxy & Spearman $\rho$ & Kendall $\tau$ & 95\% CI of $\tau$ \\
\midrule
\multicolumn{4}{l}{\textit{EndorseRank}} \\
In-degree & 0.302 & 0.218 & [0.199, 0.238] \\
In-value (sum inbound) & 0.256 & 0.175 & [0.157, 0.194] \\
\addlinespace
\multicolumn{4}{l}{\textit{AWP}} \\
In-degree & 0.883 & 0.742 & [0.732, 0.752] \\
In-value (sum inbound) & 0.887 & 0.730 & [0.720, 0.740] \\
\bottomrule
\end{tabular}
}
\end{table}


AWP agrees with the in-value ($0.566$, $[0.558,0.575]$) almost as closely as with the in-degree ($0.578$, $[0.569,0.587]$), for the reason given in Section~\ref{sec:same-window-alignment}. EndorseRank's transfer coefficients ($0.182$, $[0.169,0.194]$, and $0.134$, $[0.123,0.145]$) have intervals that overlap neither zero nor AWP's intervals. They mainly reflect that contracts that many owners approve receive transfers from more senders than the rest of the cohort. A protocol that equates reputation with inbound flow could read that flow from the in-degree itself, but out of window AWP ranks new senders above the in-degree in both windows (Section~\ref{sec:holdout-results}), so on the transfer side the walk adds to the count it smooths.

\dissertationSubsubheading[sub:allowance-family]{Allowance family}

Table~\ref{tab:alignment-allowance} gives the coefficients of the two scores with the in-approve degree and the in-approve value.

% Generated by experiments/scripts/export_latex.py -- do not edit by hand.
% Source: experiments/data/2-processed-tables-and-evaluations/revised-usd-12m/same-window/eval_summary.json
\begin{table}[htbp]
\centering
\caption[Domain alignment of EndorseRank and AWP with the allowance proxies.]{Domain alignment of EndorseRank and AWP with the allowance proxies (in-approve degree and value of each spender contract), same window (matched cohort, $n=6{,}483$ contracts). Kendall $\tau$ and Spearman $\rho$ between scores and proxy rankings. CI = 95\% percentile bootstrap over contracts (400 paired resamples, seed 42).}
\label{tab:alignment-allowance}
\fitwidth{%
\begin{tabular}{lccc}
\toprule
Proxy & Spearman $\rho$ & Kendall $\tau$ & 95\% CI of $\tau$ \\
\midrule
\multicolumn{4}{l}{\textit{EndorseRank}} \\
In-approve degree & 0.804 & 0.650 & [0.639, 0.661] \\
In-approve value (sum) & 0.628 & 0.469 & [0.454, 0.484] \\
\addlinespace
\multicolumn{4}{l}{\textit{AWP}} \\
In-approve degree & 0.320 & 0.242 & [0.223, 0.262] \\
In-approve value (sum) & 0.294 & 0.211 & [0.193, 0.229] \\
\bottomrule
\end{tabular}
}
\end{table}


EndorseRank follows the in-approve degree ($0.655$, $[0.648,0.662]$) more closely than the in-approve value ($0.515$, $[0.506,0.524]$), for the reason given in Section~\ref{sec:same-window-alignment}. AWP's coefficients are positive ($0.206$, $[0.193,0.217]$, and $0.180$, $[0.169,0.191]$): every cohort contract received approvals of the five tokens from at least three owners, and among them the contracts that many owners approve also tend to be central in the transfer graph. The comparison is partly built in, since EndorseRank is computed from the same approvals.

\dissertationSubsubheading[sub:sybil-family]{Sybil-adjusted stability family}

Table~\ref{tab:alignment-sybil-stability} gives the coefficients of the two scores with the inbound counterparty ratio, the transfer tenure and the active months.

% Generated by experiments/scripts/export_latex.py -- do not edit by hand.
% Source: experiments/data/2-processed-tables-and-evaluations/revised-usd-12m/same-window/eval_summary.json
\begin{table}[htbp]
\centering
\caption[Domain alignment of EndorseRank and AWP with the Sybil-adjusted stability proxies.]{Domain alignment of EndorseRank and AWP with the Sybil-adjusted stability proxies from transfer activity, same window (matched cohort, $n=6{,}483$ contracts). Kendall $\tau$ and Spearman $\rho$ between scores and proxy rankings. CI = 95\% percentile bootstrap over contracts (400 paired resamples, seed 42).}
\label{tab:alignment-sybil-stability}
\fitwidth{%
\begin{tabular}{lccc}
\toprule
Proxy & Spearman $\rho$ & Kendall $\tau$ & 95\% CI of $\tau$ \\
\midrule
\multicolumn{4}{l}{\textit{EndorseRank}} \\
Inbound counterparty ratio & -0.290 & -0.197 & [-0.213, -0.180] \\
Transfer tenure (days) & 0.199 & 0.133 & [0.114, 0.153] \\
Active months & 0.203 & 0.143 & [0.124, 0.164] \\
\addlinespace
\multicolumn{4}{l}{\textit{AWP}} \\
Inbound counterparty ratio & 0.163 & 0.055 & [0.033, 0.079] \\
Transfer tenure (days) & 0.763 & 0.588 & [0.575, 0.600] \\
Active months & 0.768 & 0.616 & [0.603, 0.629] \\
\bottomrule
\end{tabular}
}
\end{table}


AWP is ahead of EndorseRank on tenure ($0.508$, $[0.499,0.517]$, against $0.080$, $[0.069,0.091]$) and active months ($0.549$, $[0.540,0.558]$, against $0.098$, $[0.087,0.109]$), as one would expect from a transfer score that rewards activity spread over the window; the paired difference on the family mean excludes zero (Table~\ref{tab:tau-diff}, contrast~iv). On the inbound counterparty ratio the two have opposite signs, $-0.156$ ($[-0.166,-0.147]$) for EndorseRank and $0.127$ ($[0.114,0.141]$) for AWP. The ratio falls when the same senders return, so contracts that many owners approve are, on this measure, contracts whose senders return. AWP's positive coefficient comes from coverage: the $2{,}670$ contracts outside the transfer graph have ratio zero and AWP score zero, so they fill the bottom of both orderings. The family describes how much longitudinal transfer information each score carries.

\dissertationSubheading[sec:er-awp-holdout]{Out of window}

Table~\ref{tab:endorserank-awp-holdout} collects the paired differences between the two scores in both holdout windows. In every row $a$ is EndorseRank and $b$ AWP.

% Generated by experiments/scripts/export_latex.py -- do not edit by hand.
% Source: experiments/data/2-processed-tables-and-evaluations/revised-usd-12m/holdout-w0/eval_summary.json
% Source: experiments/data/2-processed-tables-and-evaluations/revised-usd-12m/supplementary/supplementary.json
% Source: experiments/data/2-processed-tables-and-evaluations/revised-usd-12m/window-w1/eval_summary.json
\begin{table}[htbp]
\centering
\caption[Paired bootstrap differences between EndorseRank and AWP out of window.]{Paired bootstrap differences between EndorseRank and AWP out of window (W0 spenders $n=7{,}914$; W0 traders $n=87$; W1 spenders $n=7{,}960$; W1 traders $n=70$). None of them enters rule A or rule B: the spender rows were computed after the labels were known, and the trader rows are the EndorseRank-minus-AWP pair of the comparison on an outcome built from neither edge type (Tables~\ref{tab:neutral-label-contrasts} and~\ref{tab:neutral-label-grid}). The W1 spender rows were computed from the saved W1 scores after the W1 evaluation, with the resamples of Table~\ref{tab:fresh-holdout}. Within each block $a$ and $b$ share the contract resamples (400 for the spenders and 2{,}000 for the traders; seed 42).}
\label{tab:endorserank-awp-holdout}
\fitwidth{%
\begin{tabular}{llccccc}
\toprule
Contrast ($a$ minus $b$) & Label & $\tau(a)$ & $\tau(b)$ & $\Delta\tau$ & 95\% CI of $\Delta\tau$ & $P(\Delta\tau>0)$ \\
\midrule
\multicolumn{7}{l}{\emph{W0, spenders ($n=7{,}914$)}} \\
EndorseRank minus AWP & new approval pairs & 0.303 & 0.196 & +0.107 & [0.086, 0.130] & 1.000 \\
EndorseRank minus AWP & new transfer senders & 0.272 & 0.409 & -0.138 & [-0.151, -0.123] & 0.000 \\
\midrule
\multicolumn{7}{l}{\emph{W0, traders with at least three closes ($n=87$)}} \\
EndorseRank minus AWP & liquidation-free close rate & 0.236 & 0.254 & -0.019 & [-0.173, 0.144] & 0.401 \\
\midrule
\multicolumn{7}{l}{\emph{W1, spenders ($n=7{,}960$)}} \\
EndorseRank minus AWP & new approval pairs & 0.299 & 0.157 & +0.142 & [0.119, 0.164] & 1.000 \\
EndorseRank minus AWP & new transfer senders & 0.259 & 0.398 & -0.139 & [-0.154, -0.125] & 0.000 \\
\midrule
\multicolumn{7}{l}{\emph{W1, traders with at least three closes ($n=70$)}} \\
EndorseRank minus AWP & liquidation-free close rate & 0.297 & 0.263 & +0.034 & [-0.067, 0.134] & 0.750 \\
\bottomrule
\end{tabular}
}
\end{table}


On the spenders of W0, EndorseRank ranks new approval pairs at $\tau=0.231$ (interval $[0.220,0.242]$) against $0.196$ ($[0.181,0.210]$) for AWP, a paired difference of $+0.035$ ($[0.019,0.051]$). On new transfer senders the order reverses: AWP reaches $0.346$ ($[0.336,0.356]$) against $0.200$ ($[0.189,0.210]$) for EndorseRank, a difference of $-0.146$ ($[-0.156,-0.136]$). Every label falls after all scored events, so on this cohort each edge type carries information about its own construct that the other lacks. Coverage accounts for part of the first gap. AWP scores zero for the $2{,}891$ spenders without a transfer edge before the freeze, and these gained new approval pairs more often than the rest. On the $14{,}094$ spenders that had received a transfer from another address, the order on new approval pairs reverses, AWP at $0.298$ ($[0.287,0.308]$) and EndorseRank at $0.236$ ($[0.224,0.249]$), without a paired test (Table~\ref{tab:holdout-receiving}), and AWP stays ahead on new senders.

In July--September, with the contrasts computed from the W1 scores after all labels were known, the pattern repeats. EndorseRank ranks new approval pairs at $0.234$ ($[0.222,0.244]$) against $0.169$ ($[0.157,0.183]$) for AWP ($+0.065$, $[0.051,0.081]$), and AWP ranks new senders at $0.330$ ($[0.321,0.340]$) against $0.194$ ($[0.184,0.205]$) ($-0.136$, $[-0.146,-0.127]$). On the liquidation-free close rate of the traders the two scores cannot be told apart in either window ($-0.017$, $[-0.156,0.128]$, in W0 and $-0.032$, $[-0.133,0.066]$, in W1); the trader cohorts hold $88$ and $70$ contracts.

On the W0 traders (Table~\ref{tab:holdout-traders}), AWP ranks the number of profitable closes at $0.155$ ($[0.005,0.297]$) against $0.150$ ($[-0.019,0.295]$) for EndorseRank, and the realized gain at $0.196$ ($[0.048,0.351]$) against $0.078$ ($[-0.094,0.226]$); both labels grow with the number of closes. On the two shares, which do not, both scores are negative, and for both the interval lies below zero on the non-loss share only: AWP at $-0.254$ ($[-0.404,-0.104]$) and EndorseRank at $-0.178$ ($[-0.298,-0.042]$), against $-0.124$ ($[-0.265,0.015]$) and $-0.055$ ($[-0.196,0.084]$) on the profitable share. The transfer in-degree is negative on both shares as well. Neither score ranks these contract traders by the share of successful closes.

\dissertationSubheading[sec:er-awp-neutral-grid]{All trader labels in both windows}

Table~\ref{tab:neutral-label-grid} sets the allowance side against the transfer side on all six trader labels in both windows. Its pairs were registered with W1 before the study extracted the logs of July--September (Section~\ref{sub:neutral-label-protocol}); here they are computed on the revised scores, after the registered results were known, so the table is a re-analysis (Appendix~\ref{sec:map-registered}). P and Q on the liquidation-free close rate enter rules R1 and R2, and P on the two shares enters rule R3 (Claim~C6, Section~\ref{sub:neutral-label-results}); the other entries are reported beside them. In W1 no pair separates the two sides on either liquidation label. In W0 the counts (P) favor the allowance side on both liquidation labels, while the three PageRank pairs lie below zero with intervals that include zero. On the four profit labels of W1 every pair favors the transfer side: the intervals of Q and of the EndorseRank pair lie below zero on all four, those of P on the two shares, and those of the C-PR layers on realized gain and on both shares. In W0 every interval on the profit labels includes zero except that of the C-PR layers on realized gain, a label that grows with the number of closes.

% Generated by experiments/scripts/export_latex.py -- do not edit by hand.
% Source: experiments/data/2-processed-tables-and-evaluations/revised-usd-12m/window-w1/eval_summary.json
\begin{table}[htbp]
\centering
\caption[Allowance-side minus transfer-side scores on every trader label in both windows.]{Allowance-side minus transfer-side scores, $\Delta\tau$ with 95\% paired bootstrap interval, on all six trader labels of the window W1 ($n=70$ GMX~V2 contract traders) and of W0 ($n=87$). P: in-approve degree minus transfer in-degree. Q: EndorseRank with activity restarts minus AWP (same edge weights and restart rule). C-PR: $\lambda=1$ minus $\lambda=0$ (same USD build). ER: EndorseRank minus AWP. Profitable closes and realized gain grow with the number of closes; the other four labels do not. CI = 95\% percentile bootstrap over contracts (2{,}000 paired resamples, seed 42).}
\label{tab:neutral-label-grid}
\fitwidth{%
\begin{tabular}{lcccc}
\toprule
Label & P (counts) & Q (PageRanks) & C-PR layers & ER minus AWP \\
\midrule
\multicolumn{5}{l}{\emph{Window W1, July--September 2026 ($n=70$ traders)}} \\
Liquidation-free close rate & \begin{tabular}[c]{@{}c@{}}0.141\\[-1pt]{\footnotesize [0.037, 0.240]}\end{tabular} & \begin{tabular}[c]{@{}c@{}}0.023\\[-1pt]{\footnotesize [-0.062, 0.113]}\end{tabular} & \begin{tabular}[c]{@{}c@{}}0.019\\[-1pt]{\footnotesize [-0.093, 0.128]}\end{tabular} & \begin{tabular}[c]{@{}c@{}}0.034\\[-1pt]{\footnotesize [-0.067, 0.134]}\end{tabular} \\
No liquidation & \begin{tabular}[c]{@{}c@{}}0.141\\[-1pt]{\footnotesize [0.044, 0.232]}\end{tabular} & \begin{tabular}[c]{@{}c@{}}0.027\\[-1pt]{\footnotesize [-0.066, 0.129]}\end{tabular} & \begin{tabular}[c]{@{}c@{}}0.021\\[-1pt]{\footnotesize [-0.087, 0.132]}\end{tabular} & \begin{tabular}[c]{@{}c@{}}0.038\\[-1pt]{\footnotesize [-0.069, 0.149]}\end{tabular} \\
Profitable closes & \begin{tabular}[c]{@{}c@{}}-0.115\\[-1pt]{\footnotesize [-0.257, 0.025]}\end{tabular} & \begin{tabular}[c]{@{}c@{}}-0.215\\[-1pt]{\footnotesize [-0.399, -0.042]}\end{tabular} & \begin{tabular}[c]{@{}c@{}}-0.213\\[-1pt]{\footnotesize [-0.402, -0.028]}\end{tabular} & \begin{tabular}[c]{@{}c@{}}-0.200\\[-1pt]{\footnotesize [-0.390, -0.023]}\end{tabular} \\
Realized gain & \begin{tabular}[c]{@{}c@{}}-0.109\\[-1pt]{\footnotesize [-0.247, 0.032]}\end{tabular} & \begin{tabular}[c]{@{}c@{}}-0.299\\[-1pt]{\footnotesize [-0.457, -0.146]}\end{tabular} & \begin{tabular}[c]{@{}c@{}}-0.294\\[-1pt]{\footnotesize [-0.476, -0.116]}\end{tabular} & \begin{tabular}[c]{@{}c@{}}-0.292\\[-1pt]{\footnotesize [-0.460, -0.134]}\end{tabular} \\
Profitable share & \begin{tabular}[c]{@{}c@{}}-0.162\\[-1pt]{\footnotesize [-0.295, -0.025]}\end{tabular} & \begin{tabular}[c]{@{}c@{}}-0.262\\[-1pt]{\footnotesize [-0.436, -0.100]}\end{tabular} & \begin{tabular}[c]{@{}c@{}}-0.266\\[-1pt]{\footnotesize [-0.461, -0.070]}\end{tabular} & \begin{tabular}[c]{@{}c@{}}-0.246\\[-1pt]{\footnotesize [-0.427, -0.073]}\end{tabular} \\
Non-loss share & \begin{tabular}[c]{@{}c@{}}-0.222\\[-1pt]{\footnotesize [-0.328, -0.106]}\end{tabular} & \begin{tabular}[c]{@{}c@{}}-0.324\\[-1pt]{\footnotesize [-0.471, -0.188]}\end{tabular} & \begin{tabular}[c]{@{}c@{}}-0.316\\[-1pt]{\footnotesize [-0.481, -0.152]}\end{tabular} & \begin{tabular}[c]{@{}c@{}}-0.315\\[-1pt]{\footnotesize [-0.468, -0.176]}\end{tabular} \\
\midrule
\multicolumn{5}{l}{\emph{Window W0, April--June 2026 ($n=87$ traders)}} \\
Liquidation-free close rate & \begin{tabular}[c]{@{}c@{}}0.095\\[-1pt]{\footnotesize [-0.027, 0.234]}\end{tabular} & \begin{tabular}[c]{@{}c@{}}-0.012\\[-1pt]{\footnotesize [-0.153, 0.149]}\end{tabular} & \begin{tabular}[c]{@{}c@{}}-0.057\\[-1pt]{\footnotesize [-0.209, 0.092]}\end{tabular} & \begin{tabular}[c]{@{}c@{}}-0.019\\[-1pt]{\footnotesize [-0.173, 0.144]}\end{tabular} \\
No liquidation & \begin{tabular}[c]{@{}c@{}}0.109\\[-1pt]{\footnotesize [-0.022, 0.257]}\end{tabular} & \begin{tabular}[c]{@{}c@{}}-0.005\\[-1pt]{\footnotesize [-0.162, 0.165]}\end{tabular} & \begin{tabular}[c]{@{}c@{}}-0.056\\[-1pt]{\footnotesize [-0.211, 0.098]}\end{tabular} & \begin{tabular}[c]{@{}c@{}}-0.007\\[-1pt]{\footnotesize [-0.170, 0.165]}\end{tabular} \\
Profitable closes & \begin{tabular}[c]{@{}c@{}}-0.002\\[-1pt]{\footnotesize [-0.141, 0.129]}\end{tabular} & \begin{tabular}[c]{@{}c@{}}-0.027\\[-1pt]{\footnotesize [-0.170, 0.113]}\end{tabular} & \begin{tabular}[c]{@{}c@{}}-0.089\\[-1pt]{\footnotesize [-0.244, 0.070]}\end{tabular} & \begin{tabular}[c]{@{}c@{}}-0.046\\[-1pt]{\footnotesize [-0.205, 0.113]}\end{tabular} \\
Realized gain & \begin{tabular}[c]{@{}c@{}}-0.026\\[-1pt]{\footnotesize [-0.176, 0.112]}\end{tabular} & \begin{tabular}[c]{@{}c@{}}-0.110\\[-1pt]{\footnotesize [-0.270, 0.040]}\end{tabular} & \begin{tabular}[c]{@{}c@{}}-0.257\\[-1pt]{\footnotesize [-0.447, -0.071]}\end{tabular} & \begin{tabular}[c]{@{}c@{}}-0.139\\[-1pt]{\footnotesize [-0.311, 0.022]}\end{tabular} \\
Profitable share & \begin{tabular}[c]{@{}c@{}}-0.043\\[-1pt]{\footnotesize [-0.182, 0.081]}\end{tabular} & \begin{tabular}[c]{@{}c@{}}0.044\\[-1pt]{\footnotesize [-0.084, 0.183]}\end{tabular} & \begin{tabular}[c]{@{}c@{}}0.012\\[-1pt]{\footnotesize [-0.125, 0.150]}\end{tabular} & \begin{tabular}[c]{@{}c@{}}0.074\\[-1pt]{\footnotesize [-0.066, 0.223]}\end{tabular} \\
Non-loss share & \begin{tabular}[c]{@{}c@{}}-0.074\\[-1pt]{\footnotesize [-0.208, 0.052]}\end{tabular} & \begin{tabular}[c]{@{}c@{}}0.078\\[-1pt]{\footnotesize [-0.055, 0.217]}\end{tabular} & \begin{tabular}[c]{@{}c@{}}0.047\\[-1pt]{\footnotesize [-0.090, 0.191]}\end{tabular} & \begin{tabular}[c]{@{}c@{}}0.108\\[-1pt]{\footnotesize [-0.034, 0.256]}\end{tabular} \\
\bottomrule
\end{tabular}
}
\end{table}

